\documentclass[11pt,a4paper]{article}

% =========================================================================
% PACKAGES
% =========================================================================

% Mathematics
\usepackage{amsmath,amssymb,amsthm}
\usepackage{mathtools}

% Formatting
\usepackage[margin=1in]{geometry}
\usepackage{hyperref}
\usepackage{cleveref}
\usepackage{enumitem}

% Bibliography
\usepackage[style=numeric,sorting=none,backend=bibtex]{biblatex}
\addbibresource{references.bib}

% =========================================================================
% THEOREM ENVIRONMENTS
% =========================================================================

\theoremstyle{plain}
\newtheorem{theorem}{Theorem}[section]
\newtheorem{lemma}[theorem]{Lemma}
\newtheorem{proposition}[theorem]{Proposition}
\newtheorem{corollary}[theorem]{Corollary}

\theoremstyle{definition}
\newtheorem{definition}[theorem]{Definition}
\newtheorem{example}[theorem]{Example}
\newtheorem{remark}[theorem]{Remark}
\newtheorem{conjecture}[theorem]{Conjecture}

% =========================================================================
% CUSTOM COMMANDS
% =========================================================================

\DeclareMathOperator{\tr}{tr}
\DeclareMathOperator{\diag}{diag}
\DeclareMathOperator{\Exp}{\mathbb{E}}
\newcommand{\R}{\mathbb{R}}
\newcommand{\Sn}{S^{n-1}}
\newcommand{\SigmaQC}{\Sigma_{\mathrm{QC}}}
\newcommand{\alphav}{\alpha(v)}
\newcommand{\alphac}{\alpha_{\mathrm{c}}}
\newcommand{\alphaopt}{\alpha_{\mathrm{opt}}}
\newcommand{\lammin}{\lambda_{\min}}
\newcommand{\lammax}{\lambda_{\max}}
\newcommand{\lamc}{\lambda_{\mathrm{c}}}
\newcommand{\Ocal}{\mathcal{O}}

% =========================================================================
% METADATA
% =========================================================================

\title{\textbf{The Anisotropic Quantum-Classical Boundary:} \\[0.3em]
\large Directional Regime Parameter from Observer Information Geometry}

\author{
Max Zhuravlev\thanks{Independent researcher. Email: \texttt{max@vibecodium.ai}} \\
\small \textit{Cosmological Unification Program}
}

\date{\today}

% =========================================================================
% DOCUMENT
% =========================================================================

\begin{document}

\maketitle

% =========================================================================
% ABSTRACT
% =========================================================================

\begin{abstract}
The quantum-classical boundary is standardly treated as a global property
of an observer: either quantum or classical.  Motivated by the observation
that Fisher information matrices of physical observers are generically
anisotropic---different directions in parameter space carry very different
amounts of Fisher information---we ask whether the quantum-classical
distinction itself should be direction-dependent.
In Vanchurin's neural network cosmology, the scalar regime parameter
$\alpha \in [0,1)$ classifies observers as quantum or classical.
We extend $\alpha$ to a direction-dependent function $\alpha(v)$ on
the unit sphere in the observer's parameter space, defined through
the spectral decomposition of the Fisher information matrix $F$,
and derive all topological properties of the resulting boundary
via eigenvalue analysis and differential geometry of the unit sphere.
The key result: within Vanchurin's Type~II framework,
a single observer can be simultaneously quantum
along some Fisher eigendirections and classical along others,
with the quantum-classical boundary forming a hypersurface
$\SigmaQC \subset \Sn$ whose topology is determined by the
Fisher eigenvalue spectrum.  When the Fisher matrix has $n_+$
eigenvalues above and $n_-$ below a critical threshold,
$\SigmaQC$ is homeomorphic to $S^{n_+-1} \times S^{n_--1}$.
The directional alpha is uniform (direction-independent)
if and only if the Fisher matrix is a scalar multiple of the
identity---spectral purity.  These results have no analog
in standard decoherence theory or existing information-geometric
approaches to the quantum-classical transition.
\end{abstract}

\noindent\textbf{Keywords:} Directional regime parameter, quantum-classical
boundary, Fisher information, information geometry, observer anisotropy,
Vanchurin cosmology

\vspace{1em}
\noindent\textbf{arXiv category:} quant-ph (cross-list: math-ph)

% =========================================================================
% 1. INTRODUCTION
% =========================================================================

\section{Introduction}
\label{sec:intro}

When does a physical observer behave quantum-mechanically, and when
classically?  Standard decoherence theory~\cite{zurek2003decoherence,
joos1985emergence, schlosshauer2005decoherence} answers this through
environment-induced superselection: pointer states survive interaction
with the environment, and the off-diagonal elements of the density
matrix decay at rates determined by the system-environment coupling.
The quantum-classical transition is treated as a property of the system
\emph{as a whole}---the entire density matrix decoheres on a
single timescale.

In Vanchurin's neural network cosmology~\cite{vanchurin2020world,
vanchurin2025covariant}, a different mechanism operates.  The Type~II
framework defines a metric on the observer's parameter space:
\begin{equation}
\label{eq:type2metric}
g_{\mu\nu} = M_{\mu\nu} + \beta F_{\mu\nu},
\end{equation}
where $M_{\mu\nu}$ is a mass tensor encoding structural inertia,
$F_{\mu\nu}$ is the Fisher information matrix, and
$\beta = \alpha^2/(1-\alpha)$ with $\alpha \in [0,1)$ the
\emph{regime parameter}.  When $\alpha \approx 0$, the mass tensor
dominates (classical regime); when $\alpha \approx 1$, Fisher
information dominates (quantum regime).

A companion paper~\cite{companion_amari} (Paper~\#3 of this program)
verifies that persistent observers in causally invariant hypergraph
substrates satisfy Good Regulator conditions, forcing natural gradient
learning via Amari's uniqueness theorem~\cite{amari1998natural}.
Under the ansatz $M = F^2$ for exponential family observers,
Paper~\#3 derives a scalar $\alpha$ determined by the Fisher
eigenvalue spectrum, with a quantum-classical threshold at
condition number $\kappa(F) = 2$.

The present work addresses a fundamental limitation of the scalar
treatment.  The Fisher information matrix $F$ is generically
\emph{anisotropic}: its eigenvalues $\lambda_1, \ldots, \lambda_n$
are not all equal.  Different directions in parameter space carry
different amounts of information.  This simple observation has a
profound consequence: \emph{a single observer can be simultaneously
quantum along some parameter directions and classical along others}.

We formalize this through the \emph{directional regime parameter}
$\alpha(v)$, a function on the unit sphere $\Sn$ in parameter space.
The quantum-classical boundary is not a single threshold but a
\emph{hypersurface} $\SigmaQC \subset \Sn$.  We prove that:
\begin{enumerate}[noitemsep]
\item $\alpha(v)$ depends only on direction, not on scale
      (\Cref{thm:scale_invariance}).
\item The range of $\alpha(v)$ is determined by the spread of Fisher
      eigenvalues (\Cref{thm:alpha_range}).
\item The boundary $\SigmaQC$ is a smooth submanifold whose topology
      is completely determined by the eigenvalue spectrum
      (\Cref{thm:sigma_topology}).
\item $\alpha(v)$ is direction-independent if and only if the Fisher
      matrix has spectral purity: $F = \lambda_0 I$
      (\Cref{thm:uniform_alpha}).
\end{enumerate}

\textbf{Novelty.}  The concept of a direction-dependent regime
parameter has no analog in:
\begin{itemize}[noitemsep]
\item Standard decoherence theory~\cite{zurek2003decoherence}, where
      the system decoheres as a whole;
\item Quantum Darwinism~\cite{zurek2009quantum}, where classicality
      is measured by information redundancy without directional
      structure;
\item Information-geometric approaches to quantum mechanics
      (Amari~\cite{amari1985differential, amari2016information},
      Cencov~\cite{cencov1982statistical}), which study the geometry
      of statistical manifolds but do not define a directional
      quantum-classical boundary;
\item The Free Energy Principle~\cite{friston2010free}, which
      treats prediction error minimization without directional
      decomposition of the information metric.
\end{itemize}

\textbf{Scope.}  This is a theoretical framework within Vanchurin's
Type~II cosmology.  All results are conditional on the $M = F^2$
ansatz for exponential family observers (see \Cref{sec:setup}).
We do not claim experimental predictions; we identify what
numerical experiments would distinguish this framework from
standard decoherence (\Cref{sec:predictions}).

\textbf{Organization.}  \Cref{sec:setup} reviews the Fisher matrix
and Type~II metric.  \Cref{sec:directional} defines $\alpha(v)$
and proves its main properties.  \Cref{sec:physical} gives the
physical interpretation.  \Cref{sec:spectral} connects uniform
alpha to spectral purity.  \Cref{sec:predictions} discusses
testable predictions.  \Cref{sec:discussion} concludes.


% =========================================================================
% 2. SETUP
% =========================================================================

\section{Setup: Fisher Matrix and Type II Metric}
\label{sec:setup}

\subsection{Setting}

We work within the framework of Paper~\#3~\cite{companion_amari}.
An observer $\Ocal$ in a causally invariant hypergraph
substrate~\cite{wolfram2020ruliad} maintains an internal model
$M_\theta$ parameterized by $\theta \in \Theta \subset \R^n$.
By the Good Regulator verification (Paper~\#3, Good Regulator Proposition),
the observer minimizes prediction error at its boundary, which
induces a Fisher information matrix on the parameter space:
\begin{equation}
\label{eq:fisher_def}
F_{ij}(\theta) = \Exp\!\left[
  \frac{\partial \log p_\theta}{\partial \theta^i}\,
  \frac{\partial \log p_\theta}{\partial \theta^j}
\right].
\end{equation}

We assume throughout that $F$ is positive definite (the observer has
no redundant parameters).  The spectral decomposition is:
\begin{equation}
\label{eq:spectral}
F = U \Lambda U^\top, \qquad
\Lambda = \diag(\lambda_1, \ldots, \lambda_n),
\end{equation}
with eigenvalues $0 < \lambda_1 \leq \lambda_2 \leq \cdots \leq
\lambda_n$ and orthonormal eigenvectors $\{e_1, \ldots, e_n\}$
(columns of $U$).  The condition number is
$\kappa(F) = \lambda_n / \lambda_1$.

\subsection{The \texorpdfstring{$M = F^2$}{M = F-squared} Ansatz}

Under the ansatz that the mass tensor satisfies $M = F^2$ for
exponential family observers (Paper~\#3, Mass Tensor Remark), the Type~II
metric~\eqref{eq:type2metric} becomes:
\begin{equation}
\label{eq:gmetric}
g(\alpha) = F^2 + \beta\, F, \qquad
\beta = \frac{\alpha^2}{1 - \alpha},
\end{equation}
with eigenvalues $\mu_k = \lambda_k(\lambda_k + \beta)$ for
$k = 1, \ldots, n$.

\begin{remark}[Status of the ansatz]
\label{rem:ansatz}
The $M = F^2$ relation is an \emph{ansatz} motivated by the
structure of exponential family partition functions, not an
independently verified empirical claim.  All results in this
paper are conditional on this ansatz.  Independent numerical
verification of $M = F^2$ against direct computation of the
mass tensor remains open.  See Paper~\#3,
Mass Tensor Remark, for discussion.
\end{remark}

% =========================================================================
% 3. DIRECTIONAL ALPHA
% =========================================================================

\section{The Directional Regime Parameter}
\label{sec:directional}

\subsection{Definition}

\begin{definition}[Directional regime parameter]
\label{def:alpha_v}
For a unit vector $v \in \R^n$ with $v^\top F\, v > 0$, define the
\emph{directional regime parameter}:
\begin{equation}
\label{eq:alpha_v}
\alpha(v) = \frac{v^\top F^2 v}{v^\top F^2 v + \beta\, v^\top F\, v\,}.
\end{equation}
This quantifies the relative contribution of the mass tensor
$M = F^2$ versus the Fisher tensor $F$ to the Type~II metric
along direction~$v$.
\end{definition}

\begin{remark}[Convention warning]
\label{rem:convention}
The directional $\alpha(v)$ measures the \emph{mass fraction} of
the metric along direction~$v$: large $\alpha(v)$ means the mass
tensor dominates (classical), small $\alpha(v)$ means Fisher
dominates (quantum).  This is the \textbf{opposite convention} from
Vanchurin's scalar~$\alpha$, where $\alpha \approx 1$ means
$\beta \to \infty$ (Fisher dominates, quantum regime).  The two
quantities are related but distinct: the scalar~$\alpha$ determines
the global coupling strength~$\beta$; the directional~$\alpha(v)$
resolves how that coupling manifests along each direction.
When the Fisher matrix is isotropic ($F \propto I$), $\alpha(v)$
reduces to a direction-independent constant, recovering the scalar theory.
\end{remark}

\subsection{Scale Invariance}

\begin{theorem}[Scale invariance]
\label{thm:scale_invariance}
The directional regime parameter $\alpha(v)$ depends only on the
direction of $v$, not on its magnitude: for any $c > 0$,
$\alpha(cv) = \alpha(v)$.
\end{theorem}

\begin{proof}
Both the numerator and denominator of~\eqref{eq:alpha_v} are
quadratic forms in $v$.  Replacing $v$ by $cv$ multiplies both
by $c^2$, which cancels:
\[
\alpha(cv) = \frac{c^2\, v^\top F^2 v}
                  {c^2\, v^\top F^2 v + \beta\, c^2\, v^\top F\, v}
           = \alpha(v). \qedhere
\]
\end{proof}

By \Cref{thm:scale_invariance}, $\alpha$ is well-defined as a function
$\alpha \colon \Sn \to (0,1)$ on the unit sphere in parameter space.

\subsection{Values Along Eigendirections}

\begin{proposition}[Eigendirection alpha]
\label{prop:eigendirection}
For the $k$-th Fisher eigenvector $e_k$ with eigenvalue
$\lambda_k > 0$:
\begin{equation}
\label{eq:alpha_k}
\alpha_k \coloneqq \alpha(e_k)
  = \frac{\lambda_k}{\lambda_k + \beta}.
\end{equation}
\end{proposition}

\begin{proof}
Since $F\, e_k = \lambda_k e_k$, we have $F^2 e_k =
\lambda_k^2 e_k$.  Therefore $e_k^\top F^2 e_k = \lambda_k^2$
and $e_k^\top F\, e_k = \lambda_k$.  Substituting
into~\eqref{eq:alpha_v}:
\[
\alpha(e_k)
  = \frac{\lambda_k^2}{\lambda_k^2 + \beta\,\lambda_k}
  = \frac{\lambda_k}{\lambda_k + \beta}. \qedhere
\]
\end{proof}

\subsection{Range of the Directional Alpha}

\begin{theorem}[Range of $\alpha(v)$]
\label{thm:alpha_range}
The directional regime parameter satisfies:
\begin{equation}
\label{eq:alpha_bounds}
\alpha_{\min}
  = \frac{\lambda_1}{\lambda_1 + \beta}
  \;\leq\;
  \alpha(v)
  \;\leq\;
  \frac{\lambda_n}{\lambda_n + \beta}
  = \alpha_{\max}
\end{equation}
for all unit vectors $v$.  The minimum is attained at $v = e_1$
(smallest eigenvalue direction, most quantum) and the maximum at
$v = e_n$ (largest eigenvalue direction, most classical).
\end{theorem}

\begin{proof}
Write $v = \sum_k c_k e_k$ with $\sum_k c_k^2 = 1$ in the
eigenbasis of $F$.  Then:
\begin{align}
\alpha(v)
  &= \frac{\sum_k c_k^2 \lambda_k^2}
          {\sum_k c_k^2 \lambda_k^2 + \beta \sum_k c_k^2 \lambda_k}
  \label{eq:alpha_expansion} \\
  &= \frac{\sum_k c_k^2 \lambda_k^2}
          {\sum_k c_k^2 \lambda_k(\lambda_k + \beta)}.
  \nonumber
\end{align}
Define weights $w_k = c_k^2 \lambda_k / (\sum_j c_j^2 \lambda_j)$,
which satisfy $w_k \geq 0$ and $\sum_k w_k = 1$.  Then:
\[
\alpha(v) = \sum_k w_k \cdot \frac{\lambda_k}{\lambda_k + \beta}.
\]
This is a weighted average of the values
$\lambda_k/(\lambda_k + \beta)$ with non-negative weights summing
to~1.  Since $f(x) = x/(x+\beta)$ is strictly increasing for
$x > 0$ and $\beta > 0$:
\[
\frac{\lambda_1}{\lambda_1 + \beta}
\leq \alpha(v) \leq
\frac{\lambda_n}{\lambda_n + \beta}.
\]
Equality holds at the left when all weight concentrates on
$k = 1$ (i.e., $v = e_1$) and at the right when $v = e_n$.
\end{proof}

\begin{definition}[Alpha spread]
\label{def:alpha_spread}
The \emph{alpha spread} measures directional inhomogeneity:
\begin{equation}
\label{eq:spread}
\Delta\alpha
  \coloneqq \alpha_{\max} - \alpha_{\min}
  = \frac{\beta(\lambda_n - \lambda_1)}
         {(\lambda_n + \beta)(\lambda_1 + \beta)}.
\end{equation}
\end{definition}

\begin{corollary}[Spread bounds]
\label{cor:spread_bounds}
$0 \leq \Delta\alpha < 1 - 1/\kappa(F)$, with
$\Delta\alpha = 0$ if and only if $\kappa(F) = 1$
(spectral purity).
\end{corollary}

\begin{proof}
The lower bound is immediate: $\Delta\alpha = 0$ if and only if
$\lambda_n = \lambda_1$.  For the upper bound, write
$\kappa = \lambda_n/\lambda_1$ and compute from~\eqref{eq:spread}:
\[
\Delta\alpha
  = \frac{\beta\,\lambda_1(\kappa - 1)}
         {(\kappa\,\lambda_1 + \beta)(\lambda_1 + \beta)}.
\]
Expanding the inequality $\beta\,\lambda_1\,\kappa <
(\kappa\,\lambda_1 + \beta)(\lambda_1 + \beta)$ gives
$0 < \kappa\,\lambda_1^2 + \beta\,\lambda_1 + \beta^2$,
which holds for all positive quantities.  Hence
$\Delta\alpha < (\kappa-1)/\kappa = 1 - 1/\kappa$.
\end{proof}


\subsection{The Quantum-Classical Boundary Surface}

\begin{definition}[Critical Fisher eigenvalue]
\label{def:critical}
Given a threshold $\alphac \in (0,1)$ delineating quantum
($\alpha > \alphac$) from classical ($\alpha < \alphac$) regimes,
define the \emph{critical Fisher eigenvalue}:
\begin{equation}
\label{eq:lambda_c}
\lamc = \frac{\beta\,\alphac}{1 - \alphac}.
\end{equation}
An eigendirection $e_k$ is quantum if $\lambda_k > \lamc$ and
classical if $\lambda_k < \lamc$.
\end{definition}

\begin{definition}[Quantum-classical boundary surface]
\label{def:sigma_qc}
The \emph{quantum-classical boundary surface} is:
\begin{equation}
\label{eq:sigma_qc}
\SigmaQC = \bigl\{ v \in \Sn : \alpha(v) = \alphac \bigr\}.
\end{equation}
\end{definition}

\begin{theorem}[Topology of $\SigmaQC$]
\label{thm:sigma_topology}
Assume the Fisher matrix $F$ has positive definite eigenvalues and
$\alphac \in (0,1)$.
\begin{enumerate}[label=(\alph*)]
\item $\SigmaQC$ is the level set of a smooth function on $\Sn$.

\item In the eigenbasis of $F$, $\SigmaQC$ is defined by:
\begin{equation}
\label{eq:sigma_equation}
\sum_{k=1}^n c_k^2\,\lambda_k(\lambda_k - \lamc) = 0,
\end{equation}
where $v = \sum_k c_k\, e_k$ with $\sum_k c_k^2 = 1$.

\item If $\lamc$ lies strictly between $\lambda_1$ and $\lambda_n$
(the observer has both quantum and classical eigendirections),
let $n_+ = |\{k : \lambda_k > \lamc\}|$ and
$n_- = |\{k : \lambda_k < \lamc\}|$.  Then $\SigmaQC$ is a
smooth $(n-2)$-dimensional submanifold of $\Sn$, homeomorphic to
$S^{n_+-1} \times S^{n_--1}$.

\item If all $\lambda_k > \lamc$: $\SigmaQC = \emptyset$ (entirely
quantum in every direction).

\item If all $\lambda_k < \lamc$: $\SigmaQC = \emptyset$ (entirely
classical in every direction).
\end{enumerate}
\end{theorem}

\begin{proof}
\textbf{Part (a).}
The function $\alpha \colon \Sn \to (0,1)$ is a ratio of
quadratic forms with strictly positive denominator (since $F$ is
positive definite and $\beta > 0$).  Both numerator and denominator
are smooth, so $\alpha$ is smooth on $\Sn$.

\textbf{Part (b).}
Setting $\alpha(v) = \alphac$ in~\eqref{eq:alpha_expansion} and
rearranging:
\begin{align*}
\sum_k c_k^2 \lambda_k^2 \,(1 - \alphac)
  &= \alphac\,\beta \sum_k c_k^2 \lambda_k.
\end{align*}
Defining $\lamc = \beta\,\alphac/(1-\alphac)$ gives:
\[
\sum_k c_k^2 \lambda_k^2 / \lamc = \sum_k c_k^2 \lambda_k,
\]
which rearranges to~\eqref{eq:sigma_equation}.

\textbf{Part (c).}
Partition the index set into $I_+ = \{k : \lambda_k > \lamc\}$
and $I_- = \{k : \lambda_k < \lamc\}$.
Equation~\eqref{eq:sigma_equation} becomes:
\begin{equation}
\label{eq:balance}
\sum_{k \in I_+} c_k^2\,\lambda_k(\lambda_k - \lamc)
= \sum_{k \in I_-} c_k^2\,\lambda_k(\lamc - \lambda_k).
\end{equation}
Both sides are non-negative.  Introduce rescaled coordinates
$u_k = c_k \sqrt{\lambda_k\,|\lambda_k - \lamc|}$.
Then~\eqref{eq:balance} becomes $\|u_+\|^2 = \|u_-\|^2$,
which we can write as $r_+^2 = r_-^2$, i.e., $r_+ = r_-$.

The constraint $\sum_k c_k^2 = 1$ becomes an ellipsoid condition
on the $u$-coordinates: $\sum_k u_k^2 / (\lambda_k\,|\lambda_k - \lamc|)
= 1$.  For generic eigenvalues (no $\lambda_k = \lamc$), the
surface is smooth.  To see the product topology, write
$\hat{u}_+ = u_+/\|u_+\|$ and $\hat{u}_- = u_-/\|u_-\|$ for the
directions in the $I_+$- and $I_-$-subspaces.  Given any pair
$(\hat{u}_+, \hat{u}_-)$, the balance condition $\|u_+\| = \|u_-\|$
together with the ellipsoid constraint determines a unique positive
scale $R(\hat{u}_+, \hat{u}_-)$ via
$R^2 = 1/[\sum_{k \in I_+} \hat{p}_k^2/a_k^{+} +
\sum_{k \in I_-} \hat{q}_k^2/a_k^{-}]$,
where $a_k^{\pm} = \lambda_k\,|\lambda_k - \lamc|$ are positive
constants.  The map $(\hat{u}_+, \hat{u}_-) \mapsto
(R\,\hat{u}_+,\, R\,\hat{u}_-)$ is a continuous bijection from
$S^{n_+-1} \times S^{n_--1}$ onto $\SigmaQC$ with continuous
inverse (since $R$ is positive, finite, and continuous in the
direction variables).  Hence $\SigmaQC$ is homeomorphic to
$S^{n_+-1} \times S^{n_--1}$.

Its dimension is $(n_+-1) + (n_--1) = n_+ + n_- - 2 = n - 2$
(since $n_+ + n_- = n$ when no eigenvalue equals $\lamc$ exactly).

\textbf{Parts (d)--(e).}
If all $\lambda_k > \lamc$, every term in~\eqref{eq:sigma_equation}
is positive, so the sum is strictly positive for any
$v \neq 0$---no solution exists on $\Sn$.  The case
$\lambda_k < \lamc$ for all $k$ is analogous (sum strictly
negative).
\end{proof}

\begin{remark}
The boundary $\SigmaQC$ divides $\Sn$ into a ``quantum cap''
(directions where $\alpha(v) > \alphac$) and a ``classical cap''
(where $\alpha(v) < \alphac$).  For an observer with $n_+ = 1$
quantum eigendirection and $n_- = n-1$ classical eigendirections,
$\SigmaQC \cong S^0 \times S^{n-2} \cong S^{n-2} \sqcup S^{n-2}$,
a pair of antipodal $(n-2)$-spheres.  The quantum region is a
narrow cone around the single quantum eigenvector.
\end{remark}


% =========================================================================
% 4. PHYSICAL INTERPRETATION
% =========================================================================

\section{Physical Interpretation}
\label{sec:physical}

\subsection{Direction-Dependent Quantum Effects}

The directional alpha $\alpha(v)$ answers a concrete question:
\emph{how does the observer respond to perturbations along
direction $v$ in parameter space?}

The parameter $\theta$ encodes the observer's internal model of
its environment (per the Good Regulator verification of
Paper~\#3~\cite{companion_amari}).  Different components encode
different environmental features.  The Fisher eigenvectors $e_k$
are the \emph{independent information channels}:

\begin{itemize}[noitemsep]
\item \textbf{Large $\lambda_k$} (large $\alpha_k$, near~1):
  High sensitivity to environmental signals along $e_k$.
  Under $M = F^2$, the mass eigenvalue scales as $\lambda_k^2$
  while the Fisher eigenvalue scales as $\lambda_k$, so for
  large $\lambda_k$ the mass term dominates the metric:
  $\lambda_k^2 \gg \beta\,\lambda_k$.  High inertia, precise
  measurements, slow adaptation.  This is the \emph{classical}
  regime.

\item \textbf{Small $\lambda_k$} (small $\alpha_k$, near~0):
  Low sensitivity along $e_k$.  The Fisher term dominates---low
  inertia, imprecise measurements, rapid adaptation.  This is the
  \emph{quantum} regime.
\end{itemize}

A single observer is therefore analogous to a measurement
apparatus with direction-dependent sensitivity.  Along some
directions it behaves classically (high inertia, precise
readings); along others it behaves quantum-mechanically
(low inertia, uncertain readings).

\begin{example}[Triangle graph $K_3$]
\label{ex:K3}
For the $K_3$ Ising observer at coupling $J = 0.5$
(Paper~\#3, verified across 91 configurations), the Fisher
matrix has eigenvalues $\lambda_1 = 0.385$ (multiplicity~2,
differential modes) and $\lambda_2 = 1.095$ (multiplicity~1,
uniform mode).  At $\beta = 0.5$:
\begin{align*}
\alpha_1 &= 0.385/(0.385 + 0.5) = 0.435
  &&\text{(more quantum)}, \\
\alpha_2 &= 1.095/(1.095 + 0.5) = 0.687
  &&\text{(more classical)}.
\end{align*}
The observer is more classical when measuring the overall
correlation level (high sensitivity, high inertia, single
uniform mode) and more quantum along the two differential
modes (low sensitivity, low inertia).
\end{example}

\begin{example}[Complete graph $K_5$]
\label{ex:K5}
For $K_5$ at $J = 0.5$ with $\beta = 0.5$: $\kappa(F) = 21.4$,
$\Delta\alpha = 0.536$.  The observer simultaneously spans:
\begin{itemize}[noitemsep]
\item Along the most quantum direction: $\alpha_{\min} = 0.067$
  (deep quantum regime),
\item Along the most classical direction: $\alpha_{\max} = 0.603$
  (efficient learning regime).
\end{itemize}
This single observer occupies Vanchurin's quantum regime along
some directions and efficient learning regime along others.
No scalar parameter can capture this structure.
\end{example}


\subsection{Comparison with Decoherence Theory}

Standard decoherence~\cite{zurek2003decoherence, joos1985emergence}
operates through environment-induced superselection: pointer
states survive decoherence, and off-diagonal density matrix elements
decay at a rate $\Gamma \sim \gamma\,(\Delta x/\lambda_{\mathrm{dB}})^2$.
The quantum-to-classical transition concerns the system \emph{as a
whole}.

\begin{center}
\begin{tabular}{lll}
\hline
\textbf{Aspect} & \textbf{Decoherence} & \textbf{Directional $\alpha$} \\
\hline
Mechanism & Environment coupling & Information geometry \\
Preferred basis & Pointer states & Fisher eigenvectors \\
QC boundary & Decoherence rate $\Gamma$ & $\alpha(v) = \alphac$ \\
Directionality & Not explicitly directional & Explicitly directional \\
Mathematical object & Density matrix $\rho$ & Fisher metric $F$ \\
Observer role & System decoheres passively & Observer structure matters \\
\hline
\end{tabular}
\end{center}

Three key differences deserve emphasis:

\begin{enumerate}[noitemsep]
\item \textbf{Directionality.}  Decoherence theory selects a
  preferred basis (pointer states) via the interaction Hamiltonian,
  but does not assign direction-dependent decoherence rates
  within that basis.  The directional alpha assigns each direction
  its own quantum-classical character.

\item \textbf{Observer topology.}  In decoherence theory, the
  transition depends on the system-environment coupling strength.
  In the directional alpha framework, the transition depends
  additionally on the observer's \emph{internal topology}
  (connectivity of the underlying graph), which determines the
  Fisher eigenvalue spectrum.

\item \textbf{Simultaneous multi-regime.}  Decoherence produces a
  (rapid) transition: a system is either coherent or decohered.
  Directional alpha allows a single observer to be simultaneously
  in different Vanchurin regimes along different directions, with
  no analog in the density-matrix framework.
\end{enumerate}

\begin{remark}[Honest comparison]
These frameworks operate on different mathematical objects
(density matrix vs.\ Fisher metric on parameter space) and
cannot be directly equated.  In particular, Fisher eigenvectors
live in parameter space while pointer states live in Hilbert
space.  A formal mapping between the two---if one exists---has
not been established and would require showing that Fisher
eigenvectors correspond to pointer observables in an appropriate
limit.  We treat this as an open question, not a result.
\end{remark}


\subsection{Connection to Paper \#3: Scalar Alpha as Average}

The scalar regime parameter $\alpha$ of Paper~\#3 can be
understood as an effective average over directions.  Specifically,
the convergence-time optimal $\alphaopt$ from
Paper~\#3, Optimal Alpha Theorem, minimizes a global functional
$T(\alpha) = \kappa(g(\alpha)) \cdot \mu_{\max}(g(\alpha))$ that
depends on the \emph{extreme} eigenvalues of the combined metric.
The directional alpha reveals the fine structure hidden by this
optimization: even when the optimal scalar $\alpha$ places the
observer in the ``efficient learning'' regime on average, individual
eigendirections may be deep in the quantum or classical regime.

For tree-topology observers ($P_n$, $S_n$), the Fisher matrix
satisfies $F = \text{sech}^2(J) \cdot I$ (Tree Fisher Identity,
Paper~\#1~\cite{companion_lovelock}), so all eigenvalues are
equal, $\alpha(v)$ is constant, and the scalar theory is exact.
The directional framework adds content only for observers with
nontrivial topology ($\kappa(F) > 1$).


% =========================================================================
% 5. SPECTRAL PURITY AND UNIFORM ALPHA
% =========================================================================

\section{Spectral Purity and Uniform Alpha}
\label{sec:spectral}

\begin{theorem}[Uniform alpha if and only if spectral purity]
\label{thm:uniform_alpha}
The directional regime parameter $\alpha(v)$ is independent of
direction $v$ (i.e., constant on $\Sn$) if and only if $F$ has a
single distinct nonzero eigenvalue: $F = \lambda_0 I$ for some
$\lambda_0 > 0$.
\end{theorem}

\begin{proof}
$(\Longrightarrow)$ Assume $\alpha(v)$ is constant.  In
particular, for any two eigenvectors $e_i$ and $e_j$
with eigenvalues $\lambda_i, \lambda_j > 0$:
\[
\frac{\lambda_i}{\lambda_i + \beta}
  = \frac{\lambda_j}{\lambda_j + \beta}.
\]
Cross-multiplying:
$\lambda_i(\lambda_j + \beta) = \lambda_j(\lambda_i + \beta)$,
which simplifies to $\lambda_i\,\beta = \lambda_j\,\beta$.
Since $\beta > 0$, this forces $\lambda_i = \lambda_j$.
As $i, j$ were arbitrary, all eigenvalues are equal.

$(\Longleftarrow)$ If $F = \lambda_0 I$, then for any unit $v$:
$v^\top F^2 v = \lambda_0^2$ and $v^\top F\, v = \lambda_0$.
Hence:
\[
\alpha(v) = \frac{\lambda_0^2}{\lambda_0^2 + \beta\,\lambda_0}
          = \frac{\lambda_0}{\lambda_0 + \beta},
\]
independent of $v$.
\end{proof}

\begin{corollary}[Consequences of spectral purity]
\label{cor:spectral_purity}
When $F = \lambda_0 I$ (spectral purity):
\begin{enumerate}[label=(\roman*)]
\item $\alpha(v) = \lambda_0/(\lambda_0 + \beta)$ for all $v$.
\item The mass tensor satisfies $M = F^2 = \lambda_0 F$
  (proportional to $F$).
\item The deviation tensor $\Delta = M - \kappa F$ vanishes
  identically.
\item $\SigmaQC$ is either empty or all of $\Sn$---no intermediate
  boundary exists.
\item The observer is a perfect Good Regulator in the spectral
  sense: its internal structure is proportionally matched to its
  information geometry in every direction.
\end{enumerate}
\end{corollary}

\begin{proof}
Parts (i)--(ii) follow directly from $F = \lambda_0 I$ and
$M = F^2 = \lambda_0^2 I = \lambda_0 F$.  For (iii),
$\kappa = \tr(M)/\tr(F) = n\lambda_0^2/(n\lambda_0) = \lambda_0$,
so $\Delta = M - \lambda_0 F = 0$.  Part (iv) follows because
$\alpha(v)$ is constant: either $\alpha(v) > \alphac$ for all $v$
(quantum everywhere, $\SigmaQC = \emptyset$), or
$\alpha(v) < \alphac$ for all $v$ (classical everywhere,
$\SigmaQC = \emptyset$), or $\alpha(v) = \alphac$ for all $v$
($\SigmaQC = \Sn$, a non-generic case).
Part (v) restates the connection between $M \propto F$ and
perfect Good Regulator geometry established in
Paper~\#3~\cite{companion_amari}.
\end{proof}

\subsection{Connection to the Tree Fisher Identity}

Paper~\#1~\cite{companion_lovelock} proves the \emph{Tree Fisher
Identity}: for any Ising model on a tree graph $G$ with uniform
coupling $J$, the Fisher matrix is diagonal with a single
eigenvalue:
\[
F = \text{sech}^2(J) \cdot I.
\]
By \Cref{thm:uniform_alpha}, tree-topology observers automatically
have uniform directional alpha.  This is a strong structural
result: regardless of the tree's shape (path, star, binary tree,
caterpillar, etc.), the observer is informationally isotropic.

The physical meaning is clear: tree-structured observers have no
``preferred'' information direction---they are equally sensitive
in all parameter directions.  Anisotropy requires loops in the
observer topology (cycles, complete graphs), which create
correlations between different parameter directions and break the
spectral degeneracy of $F$.


\subsection{The Deviation Tensor Perspective}

The spectral purity condition can be reformulated through the
\emph{deviation tensor} introduced in
Paper~\#3~\cite{companion_amari}:

\begin{definition}[Deviation tensor]
\label{def:deviation}
Let $\kappa_{\tr} = \tr(M)/\tr(F)$.  The deviation tensor is:
\begin{equation}
\label{eq:deviation}
\Delta_{\mu\nu} = M_{\mu\nu} - \kappa_{\tr}\,F_{\mu\nu}.
\end{equation}
It is trace-free by construction: $\tr(\Delta) = 0$.
\end{definition}

\begin{theorem}[Directional alpha determines deviation sign]
\label{thm:alpha_delta}
For $M = F^2$ and the deviation tensor
$\Delta = M - \kappa_{\tr} F$, an eigendirection $e_k$ has:
\begin{itemize}[noitemsep]
\item $\delta_k > 0$ (over-massive) if and only if
  $\alpha_k > \alpha_{\mathrm{mean}}$,
\item $\delta_k < 0$ (under-massive) if and only if
  $\alpha_k < \alpha_{\mathrm{mean}}$,
\item $\delta_k = 0$ if and only if
  $\alpha_k = \alpha_{\mathrm{mean}}$,
\end{itemize}
where $\delta_k = \lambda_k(\lambda_k - \kappa_{\tr})$ are the
eigenvalues of $\Delta$ and
$\alpha_{\mathrm{mean}} = \kappa_{\tr}/(\kappa_{\tr} + \beta)$.
\end{theorem}

\begin{proof}
Under $M = F^2$, the eigenvalues of $\Delta$ are
$\delta_k = \lambda_k^2 - \kappa_{\tr}\,\lambda_k =
\lambda_k(\lambda_k - \kappa_{\tr})$.
So $\delta_k > 0$ if and only if $\lambda_k > \kappa_{\tr}$.

Meanwhile, $\alpha_k > \alpha_{\mathrm{mean}}$ means
$\lambda_k/(\lambda_k + \beta) > \kappa_{\tr}/(\kappa_{\tr}+\beta)$.
Cross-multiplying (all terms positive):
$\lambda_k(\kappa_{\tr} + \beta) > \kappa_{\tr}(\lambda_k + \beta)$,
which simplifies to $\lambda_k\,\beta > \kappa_{\tr}\,\beta$,
i.e., $\lambda_k > \kappa_{\tr}$.  This is the same condition.
\end{proof}

The deviation tensor and directional alpha are thus complementary
views of the same phenomenon: ``over-massive'' directions (excess
structural inertia relative to information content) are precisely
the ``more classical than average'' directions.  An observer with
$\Delta = 0$ has uniform $\alpha(v)$---every direction is equally
quantum or classical.


% =========================================================================
% 6. TESTABLE PREDICTIONS
% =========================================================================

\section{Testable Predictions and Falsifiability}
\label{sec:predictions}

The directional alpha framework generates specific predictions
that are in principle falsifiable through numerical experiment.
We emphasize that these are predictions about the mathematical
framework, testable through computation, not claims about
laboratory experiments.

\subsection{Computational Predictions}

The following can be tested by computing Fisher matrices for
Ising/Boltzmann observers on various graph topologies and coupling
strengths.

\begin{enumerate}[label=(P\arabic*)]
\item \textbf{Tree isotropy} (consequence of Tree Fisher Identity).
  All tree-topology observers must have $\Delta\alpha = 0$.
  This follows mathematically from the Tree Fisher Identity
  ($F = \mathrm{sech}^2(J) \cdot I$, Paper~\#1) and
  \Cref{thm:uniform_alpha}.  A counterexample would falsify
  the Tree Fisher Identity itself.

  \emph{Status:} Verified for all tree topologies in the
  91-configuration catalogue (Paper~\#3).

\item \textbf{Complete graph anisotropy.}
  Complete graph observers $K_n$ with $n \geq 3$ must have
  $\Delta\alpha > 0$, with the spread increasing with $n$.

  \emph{Status:} Verified for $K_3$, $K_4$, $K_5$ at $\beta = 0.5$
  ($\Delta\alpha = 0.252, 0.397, 0.536$ respectively).

\item \textbf{Topology determines boundary.}
  Two observers with identical total coupling strength but
  different internal topologies (e.g., $P_5$ vs.\ $K_5$, both
  with the same mean coupling $J$) must have qualitatively
  different quantum-classical boundaries: $P_5$ has
  $\SigmaQC = \emptyset$ while $K_5$ has a nontrivial $\SigmaQC$.

\item \textbf{Condition number threshold} (consequence of Paper~\#3).
  Under Model~A convergence (Paper~\#3), only observers with
  $\kappa(F) > 2$ have nonzero optimal $\alphaopt$ (and
  therefore nontrivial directional structure at optimal beta).
  This is a theorem of Paper~\#3, not a new prediction.

\item \textbf{Boundary topology.}
  For the two-eigenvalue-group case (e.g., $K_n$ with eigenvalue
  multiplicities $n_-$ and $n_+$), $\SigmaQC$ must be
  homeomorphic to $S^{n_+-1} \times S^{n_--1}$.  This can
  be verified by explicit parameterization of the boundary.
\end{enumerate}

\subsection{Distinguishing from Standard Decoherence}

To distinguish directional alpha from standard decoherence theory,
one would need to identify a physical system where:

\begin{enumerate}[label=(D\arabic*)]
\item The decoherence rate is \emph{direction-dependent} in a way
  that correlates with Fisher eigenvalue ordering (not with the
  system-environment interaction Hamiltonian).

\item The pattern of direction-dependent decoherence rates changes
  with the observer's \emph{internal topology} while holding the
  total coupling strength fixed.

\item A single system simultaneously exhibits coherent behavior
  along some measurement directions and classical behavior along
  others, with the boundary matching the predicted $\SigmaQC$.
\end{enumerate}

These are criteria for what a decisive experiment would look like,
not claims that such experiments currently exist.

\subsection{Falsifiability}

\begin{enumerate}[label=(F\arabic*)]
\item \textbf{Spectral purity violation.}  If a tree-structured
  observer is found with $\Delta\alpha \neq 0$, this falsifies
  the Tree Fisher Identity and the framework built on it.

\item \textbf{Isotropic structured observers.}  If a
  complete-graph observer is found with $\Delta\alpha = 0$
  (isotropic decoherence despite structured topology), this
  contradicts the prediction that topology determines anisotropy.

\item \textbf{Wrong eigenvalue ordering.}  If the most rapidly
  decohering direction does not correspond to the largest Fisher
  eigenvalue, this falsifies the connection between Fisher
  spectrum and quantum-classical character.

\item \textbf{Universal scalar alpha.}  If all physical observers
  are found to have direction-independent $\alpha$ regardless of
  topology, the directional framework is unnecessary.

\item \textbf{Non-positive-definite Fisher matrices.}  If
  physical observers commonly have singular Fisher matrices in
  relevant parameter directions, the framework requires
  modification.
\end{enumerate}


% =========================================================================
% 7. DISCUSSION
% =========================================================================

\section{Discussion}
\label{sec:discussion}

\subsection{Summary of Results}

We have extended Vanchurin's scalar regime parameter $\alpha$ to a
direction-dependent function $\alpha(v)$ on the unit sphere in
observer parameter space.  The main results are:

\begin{enumerate}[noitemsep]
\item \textbf{\Cref{thm:scale_invariance}:} $\alpha(v)$ depends
  only on direction (scale-invariant).

\item \textbf{\Cref{thm:alpha_range}:} The range
  $[\alpha_{\min}, \alpha_{\max}]$ is determined by the Fisher
  eigenvalue spread.

\item \textbf{\Cref{thm:sigma_topology}:} The quantum-classical
  boundary $\SigmaQC$ is a smooth $(n-2)$-dimensional submanifold
  of $\Sn$, homeomorphic to $S^{n_+-1} \times S^{n_--1}$.

\item \textbf{\Cref{thm:uniform_alpha}:} Uniform $\alpha(v)$ if
  and only if spectral purity ($F = \lambda_0 I$).

\item \textbf{\Cref{thm:alpha_delta}:} The deviation tensor sign
  matches the directional alpha ordering.
\end{enumerate}

All five results are proven (mathematical certainty, conditional
on the $M = F^2$ ansatz).

\subsection{What This Does Not Explain}

We are explicit about the boundaries of this framework:

\begin{itemize}[noitemsep]
\item We do \emph{not} derive the $M = F^2$ ansatz from first
  principles.  It is an assumption.

\item We do \emph{not} establish a formal mapping between Fisher
  eigenvectors and Zurek's pointer states.  The analogy is
  suggestive but unproven.

\item We do \emph{not} predict specific experimental outcomes.
  The predictions are computational (testable by simulation),
  not laboratory-based.

\item We do \emph{not} explain why the critical threshold
  $\alphac$ takes any particular value.  It is a parameter of
  the framework.

\item We do \emph{not} address how $\alpha(v)$ transforms under
  observer composition or coupling.  This is an open problem
  (see \Cref{sec:open}).
\end{itemize}


\subsection{Connection to Paper \#5}

Paper~\#5~\cite{companion_spectral} of this program develops
spectral gap selection for observer Fisher matrices.  The
directional alpha framework is complementary: Paper~\#5 addresses
\emph{which} eigenvalue spectrum an observer develops, while this
paper addresses the \emph{consequences} of a given spectrum for
the quantum-classical boundary.  Together, they predict that
observers under selective pressure develop specific spectral
gap structures, which in turn determine the geometry of their
quantum-classical boundaries.

\subsection{Open Problems}
\label{sec:open}

\begin{enumerate}[noitemsep]
\item \textbf{Observer composition.}
  How does $\SigmaQC$ change when two observers are coupled?
  Can coupling create new quantum directions in a composite
  system where both components are individually classical?
  (Conjectured to be possible---analogous to entanglement
  generation---but not proven.)

\item \textbf{RG flow of directional alpha.}
  As the observer coarse-grains (renormalization group flow),
  the Fisher eigenvalue spectrum changes.  Does $\SigmaQC$
  shrink, expand, or change topology under coarse-graining?

\item \textbf{Continuous-variable extension.}
  The current framework is developed for discrete (Ising/Boltzmann)
  observers.  Extension to continuous-variable models (Gaussian
  graphical models, field theories) is needed.  Preliminary
  results suggest that tree-structured Gaussian observers do
  \emph{not} have diagonal Fisher matrices, so the spectral
  purity results may be model-dependent.

\item \textbf{Fisher--pointer state correspondence.}
  Establishing (or refuting) a formal mapping between Fisher
  eigenvectors in parameter space and pointer states in Hilbert
  space would either strengthen or weaken the physical
  interpretation of this framework.

\item \textbf{Boundary phenomenology.}
  What happens dynamically at $\SigmaQC$?  Critical slowing
  down, enhanced fluctuations, and spontaneous symmetry breaking
  are conjectured but not derived.
\end{enumerate}


\subsection{Conclusion}

The quantum-classical boundary is not a single threshold.
Within Vanchurin's Type~II framework and the $M = F^2$ ansatz,
the boundary is a hypersurface $\SigmaQC$ in the observer's
parameter space whose topology is determined by the Fisher
eigenvalue spectrum.  A single observer can be simultaneously
quantum along some information channels and classical along
others.  This anisotropic structure---absent from standard
decoherence theory---arises directly from the spectral properties
of the Fisher information matrix and provides a new perspective
on the quantum-classical transition as a per-direction,
topology-dependent spectral phenomenon.

\vspace{1em}
\noindent\textbf{Acknowledgments.}  This work was developed using
Claude Code (Anthropic) for formalization and verification.
The author thanks the Wolfram Physics Project and Vitaly Vanchurin
for foundational contributions that made this framework possible.

% =========================================================================
% BIBLIOGRAPHY
% =========================================================================

\printbibliography

\end{document}
